\section{The guarded translation and its model-theoretic soundness}
\label{sec:translation}

This section defines the base translation of $\CICm$ problems into
classical untyped first-order logic and proves it sound with respect
to the standard model.  The translation follows the shape of the
CoqHammer translation \cite[\S5.2]{CzajkaKaliszyk2018} and of the
core embedding of~\cite{Czajka2016}, with one systematic design
decision, discussed in \cref{rem:polarity}: guards attached to
\emph{hypothesis} (negative) positions are always \emph{membership}
guards $\hastype(x, \trCn{A})$, while typing information emitted in
\emph{conclusion} (positive) positions may be unfolded pointwise.
The payoff of this discipline is a bidirectional truth lemma
(\cref{thm:truth}) in a model where $\hastype$ is literally set
membership.

\subsection{The target logic}

\begin{definition}[First-order signature and semantics]
\label{def:fol}
The target is classical first-order logic with equality $\approx$
over the signature $\Sigma_{\mathrm{FOL}}$ containing: a binary
function symbol $\app$ (\emph{application}, written infix and
left-associatively, $t\app u$, often just $t\,u$); a binary predicate
$\hastype$ (\emph{the guard predicate}); a unary predicate $\PR$
(\emph{truth of an encoded proposition}); an individual constant
$\hat{c}$ for every constant $c$ of the environment (including
inductive type names and constructors), constants $\lift{\SProp}$ and
$\lift{\STypeI}$, and the \emph{lifted constants} introduced in
\cref{def:lifted} below.  We use standard Tarskian semantics;
$\FM, \nu \folmodels \varphi$ denotes truth of $\varphi$ in the
structure $\FM$ under the valuation $\nu$, and $\Phi \folproves
\varphi$ denotes classical derivability.  By G\"odel completeness,
$\Phi \folproves \varphi$ iff every model of $\Phi$ satisfies
$\varphi$.  A \emph{problem} is a pair $(\Phi, \varphi)$ of a set of
sentences (axioms) and a sentence (the conjecture); an ATP
establishes $\Phi \folproves \varphi$ by refuting $\Phi \cup
\{\neg\varphi\}$.
\end{definition}

\begin{convention}
Every $\CICm$ variable is also a first-order variable.  Proof
variables never occur in translated material
(\cref{lem:proof-confinement}), so no ``proof-erasure'' constant is
needed --- a simplification over \cite{CzajkaKaliszyk2018,
Czajka2016}, where a constant $\mathsf{prf}$/$\varepsilon$ stands
for erased proofs.
\end{convention}

\subsection{Lifted constants}

The term encoding must map binders (products, abstractions, case
expressions) and propositions-in-term-position to first-order
\emph{terms}.  As in \cite{CzajkaKaliszyk2018,Czajka2016} this is
done by \emph{lifting}: the offending subterm is replaced by a new
constant applied to the encodings of its free variables, and axioms
describing the constant are emitted.  To make the translation a
well-defined function (and to model the hash-consing performed by
the implementation) we index lifted constants by what they lift.

\begin{definition}[Lifted constants; extended environment]
\label{def:lifted}
Let $\Gamma \vdash t : T$ where $t$ is not a proof, and let
$\tup{y}{:}\tup{\rho} = \mathrm{FC}(\Gamma; t)$ be the subcontext of
$\Gamma$ declaring the free variables of $t$ (in context order;
by \cref{lem:proof-confinement} none of them are proof variables).
The \emph{lifted constant} $\lift{t}_{\tup{y}:\tup{\rho}}$ is a
first-order constant indexed by the $\alpha$-equivalence class of
the pair $(\tup{y}{:}\tup{\rho},\, t)$; distinct classes yield
distinct constants.  The \emph{extended environment} $\Env^{+}$
adds, for every such pair used by the translation, the $\CICm$
definition
\[
  \ell_{(\tup{y}:\tup{\rho};\, t)} \;\coloneqq\; \lambda
  \tup{y}{:}\tup{\rho}.\, t \;:\; \Pi \tup{y}{:}\tup{\rho}.\, T .
\]
$\Env^{+}$ is admissible whenever $\Env$ is: the new constants are
definitions, interpreted by \ref{M:var} as
$\sem{\lambda\tup{y}{:}\tup{\rho}.\,t}$.  We write $\MM$ also for the
extension of the standard model to $\Env^{+}$, and we identify the
first-order constant $\lift{t}_{\tup{y}:\tup{\rho}}$ with
$\hat{\ell}_{(\tup{y}:\tup{\rho};\,t)}$.
\end{definition}

\subsection{The translation functions}

\begin{definition}[Guards]
\label{def:guards}
For a $\Gamma$-term $A$ that is not a proposition (a data type,
$\STypeI$, or kind-level), and a first-order term $u$, the
\emph{membership guard} is
\[
  \trG{\Gamma}{u}{A} \;\coloneqq\; \hastype(u, \trC{\Gamma}{A}) .
\]
For a context $\Delta = x_1{:}A_1, \ldots, x_n{:}A_n$ extending
$\Gamma$ and a formula $\varphi$, the \emph{guarded closure} is
\[
  \langle \forall \Delta \rangle_{\Gamma}\, \varphi \;\coloneqq\;
  Q_1\, \bigl( Q_2\, ( \cdots Q_n\, \varphi ) \bigr),
  \quad
  Q_i \;=\;
  \begin{cases}
    \forall x_i.\ \trG{\Gamma_i}{x_i}{A_i} \to {-} & \text{if $A_i$ not a proposition,}\\
    \trF{\Gamma_i}{A_i} \to {-} & \text{if $A_i$ a proposition,}
  \end{cases}
\]
where $\Gamma_i = \Gamma, x_1{:}A_1, \ldots, x_{i-1}{:}A_{i-1}$.
The \emph{flattened typing formula} of a first-order term $u$ at a
$\Gamma$-type $T$ (not a proposition) is defined by recursion on
$T$:
\[
  \mathrm{Flat}_\Gamma(u, \Pi x{:}A.\,B) \;=\;
     \forall x.\ \trG{\Gamma}{x}{A} \to \mathrm{Flat}_{\Gamma, x:A}(u \app x,\ B),
  \qquad
  \mathrm{Flat}_\Gamma(u, T) \;=\; \hastype(u, \trC{\Gamma}{T})
\]
for $T$ not a product.  (In a non-propositional product
$\Pi x{:}A.\,B$ of $\CICm$ the domain $A$ is never a proposition,
by the shape of $\mathcal{R}$; so the first clause needs no
proposition case.  $\mathrm{Flat}$ coincides with the guard function
$\mathsf{G}$ of \cite[\S5.2]{CzajkaKaliszyk2018}.)
\end{definition}

\begin{definition}[Translation]
\label{def:translation}
For $\Gamma$-propositions $\varphi$ the formula translation
$\trF{\Gamma}{\varphi}$, and for $\Gamma$-terms $t$ that are neither
proofs nor sorts other than $\SProp, \STypeI$ the term translation
$\trC{\Gamma}{t}$, are defined by mutual recursion.  Each clause of
$\trC{}{}$ that introduces a lifted constant also \emph{emits}
axioms into the axiom store $\Ax$; emission is idempotent (axioms
are indexed by the same $\alpha$-classes as the constants).

\medskip
\noindent\emph{Formula translation} $\trF{\Gamma}{\varphi}$:
\begin{enumerate}[label=(F\arabic*),leftmargin=3em]
\item $\trF{\Gamma}{\Pi x{:}A.\,\varphi} = \trF{\Gamma}{A} \to
  \trF{\Gamma}{\varphi}$ \quad if $A$ is a proposition
  (then $x \notin \FV(\varphi)$ by \cref{cor:impl-nondep});
\item $\trF{\Gamma}{\Pi x{:}A.\,\varphi} = \forall x.\
  \trG{\Gamma}{x}{A} \to \trF{\Gamma, x:A}{\varphi}$ \quad if $A$ is
  not a proposition;
\item $\trF{\Gamma}{\top} = \top$, \ $\trF{\Gamma}{\bot} = \bot$, \
  $\trF{\Gamma}{\varphi \circ \psi} = \trF{\Gamma}{\varphi} \circ
  \trF{\Gamma}{\psi}$ for $\circ \in \{\wedge, \vee,
  \leftrightarrow\}$;
\item $\trF{\Gamma}{t \ceq_A u} = \bigl(\trC{\Gamma}{t} \approx
  \trC{\Gamma}{u}\bigr)$;
\item $\trF{\Gamma}{\cexists{x}{A}\varphi} = \exists x.\
  \trG{\Gamma}{x}{A} \wedge \trF{\Gamma, x:A}{\varphi}$;
\item otherwise ($\varphi$ atomic: a variable, constant, application
  or case expression of type $\SProp$): $\trF{\Gamma}{\varphi} =
  \PR(\trC{\Gamma}{\varphi})$.
\end{enumerate}

\noindent\emph{Term translation} $\trC{\Gamma}{t}$:
\begin{enumerate}[label=(C\arabic*),leftmargin=3em]
\item $\trC{\Gamma}{x} = x$; \ $\trC{\Gamma}{c} = \hat{c}$; \
  $\trC{\Gamma}{\SProp} = \lift{\SProp}$; \
  $\trC{\Gamma}{\STypeI} = \lift{\STypeI}$;
\item $\trC{\Gamma}{t\,u} = \trC{\Gamma}{t} \app \trC{\Gamma}{u}$
  (by \cref{lem:proof-confinement} neither $t$ nor $u$ is a proof);
\item \emph{(proposition lifting)} if $t$ is a proposition that is
  not a variable, constant or application whose translation is
  wanted as a term (i.e.\ $t$ occurs in term position), then
  $\trC{\Gamma}{t} = \lift{t}_{\tup{y}:\tup{\rho}}\,\tup{y}$ with
  $\tup{y}{:}\tup{\rho} = \mathrm{FC}(\Gamma; t)$, emitting
  \[
    \langle \forall\, \tup{y}{:}\tup{\rho} \rangle\,
    \Bigl( \PR\bigl(\lift{t}\,\tup{y}\bigr) \leftrightarrow
           \trF{\Gamma}{t} \Bigr) ;
    \tag{$\Ax$-prop}
  \]
  for $t$ a variable, constant or application, clauses (C1)--(C2)
  apply unchanged also when $t$ is a proposition;
\item \emph{(product lifting)} if $t = \Pi x{:}A.\,B$ is not a
  proposition (so it is a data type or kind-level), then
  $\trC{\Gamma}{t} = \lift{t}_{\tup{y}:\tup{\rho}}\,\tup{y}$,
  emitting the \emph{application-typing axiom}
  \[
    \langle \forall\, \tup{y}{:}\tup{\rho} \rangle\,
    \Bigl( \forall z.\ \hastype\bigl(z, \lift{t}\,\tup{y}\bigr) \to
      \forall x.\ \trG{\Gamma}{x}{A} \to
      \hastype\bigl(z \app x,\ \trC{\Gamma,x:A}{B}\bigr) \Bigr) ;
    \tag{$\Ax$-app}
  \]
\item \emph{(abstraction lifting)} if $t = \lambda
  \tup{x}{:}\tup{A}.\,b$ with maximal $\lambda$-prefix (so $b$ is
  not an abstraction), $t$ not a proof, then $\trC{\Gamma}{t} =
  \lift{t}_{\tup{y}:\tup{\rho}}\,\tup{y}$, emitting
  \[
    \langle \forall\, \tup{y}{:}\tup{\rho},\, \tup{x}{:}\tup{A}
    \rangle\, \Bigl( \lift{t}\,\tup{y}\,\tup{x} \;\approx_{\Gamma'}\;
    b \Bigr)
    \tag{$\Ax$-$\lambda$}
  \]
  where $\Gamma' = \Gamma, \tup{x}{:}\tup{A}$ and
  $\approx_{\Gamma'}$ is the \emph{levelled equality}: for
  first-order term $u$ and $\Gamma'$-term $b$,
  \[
    u \approx_{\Gamma'} b \;\coloneqq\;
    \begin{cases}
      \PR(u) \leftrightarrow \trF{\Gamma'}{b} & \text{if $b$ is a proposition,}\\
      \forall z.\ \hastype(z, u) \leftrightarrow
        \hastype(z, \trC{\Gamma'}{b}) & \text{if $b$ is a data type or kind-level,}\\
      u \approx \trC{\Gamma'}{b} & \text{if $b$ is a data object;}
    \end{cases}
  \]
  additionally the \emph{typing atoms and flattened typing axiom}
  for $\lift{t}$ are emitted:
  \[
    \langle \forall\, \tup{y}{:}\tup{\rho} \rangle\,
      \hastype\bigl(\lift{t}\,\tup{y},\ \trC{\Gamma}{T}\bigr),
    \qquad
    \langle \forall\, \tup{y}{:}\tup{\rho} \rangle\,
      \mathrm{Flat}_{\Gamma}\bigl(\lift{t}\,\tup{y},\ T\bigr)
    \tag{$\Ax$-ty}
  \]
  where $T$ is the type of $t$ (a non-propositional product, since
  $t$ is not a proof);
\item \emph{(case lifting)} if $t = \case{I,\tup{\tau},Q}{u}{\tup{u}}$
  is not a proof and not a proposition-in-term-position handled by
  (C3), then $\trC{\Gamma}{t} =
  \lift{t}_{\tup{y}:\tup{\rho}}\,\tup{y}$, emitting the axioms
  ($\Ax$-ty) for its type $Q\,u$ --- when that type is not a product
  the flattened axiom degenerates to the typing atom --- and the
  \emph{case equation}
  \[
    \langle \forall\, \tup{y}{:}\tup{\rho} \rangle\,
    \bigvee_{j=1}^{k}\;
    \exists \tup{x}_j.\ \Bigl(
      \textstyle\bigwedge_l \trG{}{x_{j,l}}{\sigma_{j,l}[\tup{\tau}/\tup{A}]}
      \;\wedge\;
      \trC{\Gamma}{u} \approx
        \hat{c}_j\, \trC{\Gamma}{\tup{\tau}}\, \tup{x}_j
      \;\wedge\;
      \bigl( \lift{t}\,\tup{y} \approx_{\Gamma_j} u_j\,\tup{x}_j \bigr)
    \Bigr)
    \tag{$\Ax$-case}
  \]
  with $\Gamma_j = \Gamma, \tup{x}_j : \tup{\sigma}_j[\tup{\tau}/\tup{A}]$.
  Note that the guards on $\tup{y}$ --- in particular on the free
  variables of the scrutinee $u$ --- are part of the axiom; these
  are the guards that \cite[\S5.6]{CzajkaKaliszyk2018} identifies as
  never omittable, and \cref{prop:nec-mono} below reproves that
  observation semantically.
\end{enumerate}
\end{definition}

\begin{definition}[Translation of declarations and problems]
\label{def:transl-problem}
Fix an admissible environment.  For each kind of declaration the
following axioms are generated (all in the empty context).
\begin{enumerate}[label=(D\arabic*),leftmargin=3em]
\item \emph{Parameter $c : T$.}  If $T$ is a proposition: the
  \emph{premise formula} $\trFn{T}$ (labelled $c$).  Otherwise: the
  typing atom $\hastype(\hat{c}, \trCn{T})$ and the flattened axiom
  $\mathrm{Flat}_{\langle\rangle}(\hat{c}, T)$.
\item \emph{Definition $c \coloneqq t : T$.}  The axioms of (D1) for
  $c : T$, and additionally the \emph{unfolding axiom}
  $\hat{c} \approx_{\langle\rangle} t$ --- with $\approx_{\Gamma}$
  the levelled equality of ($\Ax$-$\lambda$); for $T$ kind-level of
  product shape the unfolding is emitted in the pointwise form
  $\langle\forall \tup{x}{:}\tup{A}\rangle\, (\hat{c}\,\tup{x}
  \approx_{\tup{x}:\tup{A}} t\,\tup{x})$ so that no
  $\approx$-equation between type-level terms is ever generated
  (cf.\ \cref{lem:no-type-eq}).
\item \emph{Inductive declaration} (notation of \cref{def:env}(iii);
  write $\tup{A}$ for the parameters and abbreviate $\sigma'_{j,l} =
  \sigma_{j,l}$, all in context $\tup{A}{:}\STypeI$):
  \begin{itemize}
  \item typing atoms and flattened typing axioms (D1-style) for $I$
    and for each constructor $c_j$;
  \item for each $j$, the \emph{injectivity axiom}
    $\trFn{\forall \tup{A}{:}\STypeI.\ \forall \tup{x}, \tup{x}'.\
      c_j\,\tup{A}\,\tup{x} \ceq c_j\,\tup{A}\,\tup{x}' \to
      \bigwedge_l x_l \ceq x'_l}$;
  \item for each $i \neq j$, the \emph{discrimination axiom}
    $\trFn{\forall \tup{A}{:}\STypeI.\ \forall \tup{x}, \tup{x}'.\
      \neg\,( c_i\,\tup{A}\,\tup{x} \ceq c_j\,\tup{A}\,\tup{x}' )}$;
  \item the \emph{inversion axiom}
    $\trFn{\forall \tup{A}{:}\STypeI.\ \forall z{:}I\,\tup{A}.\
      \bigvee_j \exists \tup{x}_j{:}\tup{\sigma}_j.\
      z \ceq c_j\,\tup{A}\,\tup{x}_j}$
  \end{itemize}
  (the quantified bodies are the evident $\CICm$ propositions; their
  $\mathsf{F}$-translations expand the guards).
\item \emph{Universe atom.} $\hastype(\lift{\SProp}, \lift{\STypeI})$.
  \textbf{No} axiom $\hastype(\lift{\STypeI}, \lift{\STypeI})$ is
  emitted, in contrast to \cite[\S5.4]{CzajkaKaliszyk2018}: that
  atom is false in the intended semantics, and in the stratified
  calculus $\CICm$ it is never needed.
\end{enumerate}
Given a finite set $P$ of \emph{premises} (parameter declarations
$p_i : \varphi_i$ with $\varphi_i$ closed propositions) and a closed
goal proposition $\varphi_G$, the \emph{translated problem} is
\[
  \mathrm{Transl}(P, \varphi_G) \;=\;
  \bigl( \Ax_{\mathrm{gen}} \cup \setof{\trFn{\varphi_i}}{p_i \in P},\
         \trFn{\varphi_G} \bigr)
\]
where $\Ax_{\mathrm{gen}}$ collects (D2)--(D4) axioms for all
declarations reachable from $P \cup \{\varphi_G\}$ together with all
axioms emitted by the recursive calls of $\mathsf{C}$.
\end{definition}

\begin{lemma}[No type-level equations]
\label{lem:no-type-eq}
No axiom or conjecture of a translated problem contains an
$\approx$-literal either side of which is the translation of a data
type or kind-level term.  All equations between such terms are
rendered in the extensional form $\forall z.\, \hastype(z,\cdot)
\leftrightarrow \hastype(z,\cdot)$ by the levelled equality.
\end{lemma}

\begin{proof}
By inspection of \cref{def:translation,def:transl-problem}:
$\approx$ is produced only by (F4), where both sides are data
objects by the typing rule of $\ceq$; by ($\Ax$-$\lambda$),
($\Ax$-case) and (D2) through the levelled equality, whose
$\approx$-clause is restricted to data-object bodies; and by
($\Ax$-case) in the scrutinee equation, whose sides are data
objects.
\end{proof}

\subsection{The induced structure \texorpdfstring{$\FMstar$}{M*}}

\begin{definition}[Standard first-order structure]
\label{def:mstar}
The structure $\FMstar$ has domain $\Dom \coloneqq \Utwo$, and
interprets:
\begin{itemize}
\item $u \app v$: \ $\ap(f, a) \coloneqq f(a)$ if $f$ is a function
  (graph) with $a \in \dom(f)$, and $\emptyset$ otherwise;
\item $\hastype(u, v)$: \ the membership relation
  $\setof{(a,b) \in \Dom^2}{a \in b}$;
\item $\PR(u)$: \ $\setof{a \in \Dom}{\emptyset \in a}$;
\item $\approx$: \ identity;
\item each constant $\hat{c}$ (including lifted constants, via
  \cref{def:lifted}) by $\sem{c} \in \Dom$; and
  $\lift{\SProp} \mapsto \Om$, $\lift{\STypeI} \mapsto \Uone$.
\end{itemize}
A pair $(\rho, \nu)$ of a $\Gamma$-assignment and a first-order
valuation is \emph{conforming} if $\nu(x) = \rho(x)$ for every
non-proof variable $x$ of $\Gamma$.
\end{definition}

\begin{lemma}[Compositionality]
\label{lem:compositionality}
Let $t$ be a $\Gamma$-term in the domain of $\trC{\Gamma}{-}$ and
let $(\rho, \nu)$ be conforming.  Then the value of
$\trC{\Gamma}{t}$ in $\FMstar$ under $\nu$ equals $\semr{t}{\rho}$.
\end{lemma}

\begin{proof}
Induction on $t$.  Variables and constants: by conformity and
\cref{def:mstar}; sorts: $\sem{\SProp} = \Om$, $\sem{\STypeI} =
\Uone$ by \ref{M:sorts}.  Application $t\,u$: by the induction
hypothesis the arguments of $\app$ evaluate to $\semr{t}{\rho}$ and
$\semr{u}{\rho}$.  Since $t$ is not a proof, its type is a
non-propositional product $\Pi x{:}A.\,B$, so by \ref{M:typing} and
\ref{M:pi} the value $\semr{t}{\rho}$ is a graph with domain
$\semr{A}{\rho} \ni \semr{u}{\rho}$; hence the $\ap$-clause applies
and yields $\semr{t}{\rho}(\semr{u}{\rho}) = \semr{t\,u}{\rho}$ by
\ref{M:lam}.  Lifted clauses (C3)--(C6): the translation is
$\lift{t}\,\tup{y}$, whose value is
$\ap(\cdots\ap(\sem{\ell}, \rho(y_1))\cdots, \rho(y_m))$ with $\ell
\coloneqq \lambda\tup{y}{:}\tup{\rho}.\,t$.  Each application step
is within the successive domains by \ref{M:typing} (the
$y_i$-values conform), so by \ref{M:lam}, \ref{M:conv} and
\ref{M:subst} the value equals
$\semr{(\lambda\tup{y}.\,t)\,\tup{y}}{\rho} = \semr{t}{\rho}$.
\end{proof}

\subsection{The truth lemma and base soundness}

\begin{theorem}[Truth lemma]
\label{thm:truth}
Let $\varphi$ be a $\Gamma$-proposition and $(\rho, \nu)$
conforming.  Then
\[
  \FMstar, \nu \folmodels \trF{\Gamma}{\varphi}
  \qquad\Longleftrightarrow\qquad
  \semr{\varphi}{\rho} = \holds .
\]
Moreover every axiom emitted by the translation of any term
occurring in the problem is true in $\FMstar$.
\end{theorem}

\begin{proof}
Both statements are proved simultaneously by induction on the size
of the translated term (the axioms emitted for a lifted subterm
mention only translations of terms of smaller or equal size in
extended contexts; we make the measure precise below).  Throughout,
``true'' means true in $\FMstar$ under the ambient valuation, and we
use \cref{lem:compositionality} silently for all $\mathsf{C}$-images.

\emph{Measure.}  Define the weight of a term as its size, and order
proof obligations as follows: the truth-lemma obligation for
$\varphi$ depends on truth-lemma obligations for immediate
subformulas of $\varphi$ (smaller) and on nothing else; the
axiom-truth obligations for a lifted constant
$\lift{t}_{\tup{y}:\tup{\rho}}$ depend on truth-lemma obligations
for terms of size at most that of $t$'s components ($b$, branch
bodies, $A$, $B$), each smaller than $t$.  So the induction is
well-founded.

\emph{Truth lemma, case (F1)} $\varphi = \Pi x{:}A.\,\psi$ with $A$
a proposition, $x \notin \FV(\psi)$.  By \ref{M:pi},
$\semr{\varphi}{\rho} = \holds$ iff ($\semr{A}{\rho} = \emptyset$ or
$\semr{\psi}{\rho} = \holds$): indeed $\semr{A}{\rho} \in \Om$, and
the trace-product condition ``$\forall a \in \semr{A}{\rho}.\
\semr{\psi}{\rho[x \mapsto a]} = \holds$'' is vacuous when
$\semr{A}{\rho} = \emptyset$ and equivalent to $\semr{\psi}{\rho} =
\holds$ when $\semr{A}{\rho} = \holds$ (as $x \notin \FV(\psi)$,
by \ref{M:subst} the value does not depend on $a$).  Classically,
the right-hand side is exactly the truth condition of
$\trF{}{A} \to \trF{}{\psi}$ by the induction hypothesis for $A$ and
$\psi$.

\emph{Case (F2)} $\varphi = \Pi x{:}A.\,\psi$, $A$ not a
proposition.  The translation is $\forall x.\ \hastype(x,
\trC{\Gamma}{A}) \to \trF{\Gamma,x:A}{\psi}$.  Since
$\hastype^{\FMstar}$ is membership and $\trC{\Gamma}{A}$ evaluates
to $\semr{A}{\rho}$, the translated formula is true iff for every
$d \in \semr{A}{\rho}$ the formula $\trF{\Gamma,x:A}{\psi}$ is true
under $\nu[x \mapsto d]$.  For each such $d$, the pair
$(\rho[x \mapsto d], \nu[x \mapsto d])$ is conforming for
$\Gamma, x{:}A$, and the induction hypothesis turns the condition
into: for every $d \in \semr{A}{\rho}$,
$\semr{\psi}{\rho[x \mapsto d]} = \holds$ --- which by \ref{M:pi} is
exactly $\semr{\varphi}{\rho} = \holds$.  Note this equivalence is
\emph{exact} in both directions: the guard carves out precisely the
set $\semr{A}{\rho}$ from the domain $\Dom$.  This is the point of
using membership guards in hypothesis positions.

\emph{Case (F3), (F4), (F5)}: immediate from \ref{M:conn}, the
induction hypothesis, and (for F4) \cref{lem:compositionality};
for (F5) note again that the guard in $\exists x.\ \hastype(x,
\trC{\Gamma}{A}) \wedge \cdots$ ranges exactly over
$\semr{A}{\rho}$.

\emph{Case (F6)} $\varphi$ atomic.  $\trF{\Gamma}{\varphi} =
\PR(\trC{\Gamma}{\varphi})$, and $\trC{\Gamma}{\varphi}$ evaluates
to $\semr{\varphi}{\rho} \in \Om$ (by \ref{M:typing}).  Now
$\PR^{\FMstar}(a)$ iff $\emptyset \in a$, and for $a \in \Om$ this
is equivalent to $a = \holds$.

\emph{Axiom truth.}  We verify each emitted axiom family.

($\Ax$-prop): under any values $\tup{d}$ of $\tup{y}$ passing their
guards, the pair $(\rho, \nu)$ with $\rho = [\tup{y} \mapsto
\tup{d}]$ is conforming; the left side $\PR(\lift{t}\,\tup{y})$
holds iff $\emptyset \in \semr{t}{\rho}$ iff $\semr{t}{\rho} =
\holds$, and the right side iff the same by the induction
hypothesis (truth lemma for $t$, smaller than the lifting trigger).
Values of $\tup{y}$ \emph{failing} their guards make the guarded
closure vacuously true.

($\Ax$-app): let $\tup{d}$ pass the guards, let $e$ satisfy
$\hastype(z, \lift{t}\,\tup{y})$, i.e.\ $e \in \semr{\Pi
x{:}A.\,B}{\rho}$, and let $a \in \semr{A}{\rho}$.  By \ref{M:pi},
$e$ is a graph with domain $\semr{A}{\rho}$ and $e(a) \in
\semr{B}{\rho[x \mapsto a]}$; by the $\ap$-clause, $z \app x$
evaluates to $e(a)$, and $\trC{\Gamma,x:A}{B}$ evaluates to
$\semr{B}{\rho[x\mapsto a]}$ by \cref{lem:compositionality}.  So the
conclusion atom is membership of $e(a)$ in exactly that set: true.

($\Ax$-$\lambda$): under guard-passing values, both sides of the
levelled equality evaluate compatibly: $\lift{t}\,\tup{y}\,\tup{x}$
evaluates, as in \cref{lem:compositionality}, to
$\semr{b}{\rho}$; the right side is $\trC{}{b}$ (or
$\trF{}{b}$ under $\PR$, or the $\hastype$-extension of
$\trC{}{b}$), which denotes/holds-of the same set by
\cref{lem:compositionality} (respectively by the induction
hypothesis for the $\PR$-clause; for the type-level clause both
$\hastype$-atoms test membership in the same set, so the
$\forall z$-equivalence is true).

($\Ax$-ty): the typing atom asserts $\semr{\ell\,\tup{y}}{} \in
\semr{T}{\rho}$ --- an instance of \ref{M:typing}.  The flattened
axiom: by induction on $T$'s product prefix, repeatedly applying
\ref{M:pi} exactly as in the ($\Ax$-app) case: guard-passing
arguments lie in the successive domains, applications stay in the
successive codomains, and the final atom is membership of the fully
applied value in the final codomain --- true by \ref{M:typing} and
\ref{M:pi}.

($\Ax$-case): let $\tup{d}$ pass the guards of $\tup{y}$; write
$\rho$ for the induced assignment.  By \ref{M:typing},
$\semr{u}{\rho} \in \sem{I}(\semr{\tup{\tau}}{\rho})$, so by
\ref{M:ind}(d) there are unique $j$ and arguments $\tup{e}$ (in the
appropriate denotations, \ref{M:ind}(a)) with $\semr{u}{\rho} =
\sem{c_j}(\semr{\tup{\tau}}{\rho})(\tup{e})$.  Choose the $j$-th
disjunct and witnesses $\tup{e}$: the guards on $\tup{x}_j$ hold by
construction; the scrutinee equation holds by
\cref{lem:compositionality}; and the levelled equality between
$\lift{t}\,\tup{y}$ and $u_j\,\tup{x}_j$ holds because by
\ref{M:case} $\semr{t}{\rho} = \semr{u_j}{\rho}(\tup{e}) =
\semr{u_j\,\tup{x}_j}{\rho[\tup{x}_j \mapsto \tup{e}]}$, and the
levelled equality's truth then follows exactly as in the
($\Ax$-$\lambda$) case.

(D1)/(D2): typing atoms and flattened axioms as in ($\Ax$-ty);
premise formulas are handled by \cref{thm:base-soundness} below;
unfolding axioms: by \ref{M:var} $\sem{c} = \sem{t}$, and the
levelled equality between two names of the same set is true as in
($\Ax$-$\lambda$).

(D3): each of the injectivity, discrimination and inversion axioms
is $\trFn{\chi}$ for a closed $\CICm$ proposition $\chi$; by the
truth-lemma part (already available for $\chi$, which is not larger
than the trigger --- these axioms are emitted at declaration
translation time, and we simply order them after all truth-lemma
obligations for their own subterms), it suffices that $\chi$ is
valid.  Validity of injectivity is \ref{M:ind}(b); of
discrimination, \ref{M:ind}(c); of inversion, \ref{M:ind}(d)
(every element of $\sem{I}(\tup{a})$ is a constructor value with
arguments in the required denotations).

(D4): $\Om \in \Uone$: true, as $\Om \in \Uone$ by
\cref{def:model}.
\end{proof}

\begin{theorem}[Base soundness]
\label{thm:base-soundness}
Let $\Env$ be admissible with standard model $\MM$, let $P$ be
premises whose statements $\varphi_i$ are valid in $\MM$ (e.g.\
provable, cf.\ \cref{prop:prov-valid}), and let $\varphi_G$ be a
closed goal proposition.  If
\[
  \Ax_{\mathrm{gen}} \cup \setof{\trFn{\varphi_i}}{i} \;\folproves\;
  \trFn{\varphi_G},
\]
then $\varphi_G$ is valid in $\MM$.  In particular the axiom set
$\Ax_{\mathrm{gen}} \cup \{\trFn{\varphi_i}\}_i$ is consistent.
\end{theorem}

\begin{proof}
By \cref{thm:truth}, $\FMstar$ satisfies $\Ax_{\mathrm{gen}}$, and
satisfies each $\trFn{\varphi_i}$ since $\varphi_i$ is valid.  By
soundness of classical provability, $\FMstar \folmodels
\trFn{\varphi_G}$.  By the converse direction of \cref{thm:truth},
$\sem{\varphi_G} = \holds$.  Consistency: $\FMstar$ is a model of
the axiom set.
\end{proof}

We record the invariant that the rest of the paper preserves.

\begin{definition}[Soundness invariant]
\label{def:si}
A first-order problem $(\Phi, \gamma)$ is \emph{$\MM$-sound} if
there exists an expansion $\FM$ of $\FMstar$ by interpretations of
finitely many fresh symbols such that (i) $\FM \folmodels \Phi$, and
(ii) $\FM \folmodels \gamma$ implies that $\varphi_G$ is valid in
$\MM$, where $\varphi_G$ is the pending goal.  If $(\Phi, \gamma)$
is $\MM$-sound and $\Phi \folproves \gamma$, then $\varphi_G$ is
valid.  \Cref{thm:base-soundness} states that
$\mathrm{Transl}(P, \varphi_G)$ is $\MM$-sound (with $\FM =
\FMstar$ and (ii) supplied by \cref{thm:truth}).
\end{definition}

\subsection{Discussion: design choices and the implemented translation}
\label{sec:transl-discussion}

\begin{remark}[The polarity discipline]
\label{rem:polarity}
Typing information appears in two polarities.  In \emph{hypothesis}
position --- the guard of a universally quantified premise variable
--- our translation always uses the membership guard $\hastype(x,
\trCn{A})$.  In \emph{conclusion} position --- typing axioms for
constants and lifted functions --- it emits both the membership atom
and the pointwise $\mathrm{Flat}$ axiom.  Both conclusion forms are
true in $\FMstar$ (\cref{thm:truth}); the pointwise form in
hypothesis position would \emph{not} be: the formula
$\forall f. (\forall x.\ \hastype(x, \hat{A}) \to \hastype(f \app
x, \hat{B})) \to \trFn{\psi}$ quantifies also over graphs whose
domain strictly contains $\sem{A}$; such junk elements satisfy the
pointwise hypothesis but need not satisfy $\psi$'s translation, so
the translated premise could be false in $\FMstar$ even when the
$\CICm$ premise is valid.  For the same reason a
\emph{bidirectional} characterisation axiom $\forall z.\
\hastype(z, \lift{\Pi x{:}A.B}) \leftrightarrow (\forall x.\ldots)$
--- which the CoqHammer implementation emits for lifted type
constants (\texttt{opt\_type\_lifting}) --- is unsound for
$\FMstar$: its right-to-left direction asserts that the pointwise
class is exactly $\sem{\Pi x{:}A.\,B}$, which fails in any
membership interpretation.  Our ($\Ax$-app) keeps only the sound
elimination direction, exactly as the axiom family $\Delta_{G_1}$
of \cite[Def.~45]{Czajka2016} does.  The translation
of~\cite[\S5.2]{CzajkaKaliszyk2018} uses the pointwise form
(via its function $\mathsf{G}$) also in hypothesis positions; that
choice is one of the reasons the implemented translation ``is
neither sound nor complete'' in the strict sense
\cite[\S5.6]{CzajkaKaliszyk2018}, while its guarded core in the
discipline used here is validated both proof-theoretically
\cite{Czajka2016} and model-theoretically (\cref{thm:truth}).
\end{remark}

\begin{remark}[Functional extensionality]
\label{rem:funext}
In $\FMstar$, elements of $\sem{A \to B}$ are graphs with domain
exactly $\sem{A}$, so pointwise equal elements are equal:
the translation of the functional extensionality axiom is
\emph{true} in $\FMstar$, and may soundly appear among the
premises.  The inconsistency discussed in
\cite[\S5.6]{CzajkaKaliszyk2018} --- functional extensionality
plus Coq equality mapped to $\approx$ --- arises only in
combination with the \emph{omission of guards} on the free and bound
variables of lifted abstractions, which the implementation performs
by default and our base translation does not.  Once guards are kept,
the extensionality instance forced by ($\Ax$-$\lambda$) applies only
to guard-passing arguments, and truth in $\FMstar$ is restored.
The systematic study of which guards may be dropped is exactly the
subject of \cref{sec:erasure}; the erasure criterion proved there
authorises dropping guards only at types where no collapse of the
kind exhibited in \cref{prop:nec-mono} can occur.
\end{remark}

\begin{remark}[Incompleteness]
\label{rem:incompleteness}
The translation is deliberately incomplete, and no claim of
completeness is made anywhere in this paper.  Known sources include
the absence of $\xi$/extensionality reasoning for lifted constants
(\cite[Example~51]{Czajka2016}) and the loss of higher-order
instantiation.  Completeness statements below
(\cref{prop:mono-conservative,thm:spec-faithful,prop:erasure-complete})
are always \emph{relative}: transformation $X$ does not lose
first-order proofs that existed before $X$.
\end{remark}

\begin{example}
\label{ex:translation}
Let $\Env$ contain $\mathsf{nat}$ (constructors $\mathsf{O},
\mathsf{S}$), the parametric $\mathsf{list}$ (constructors
$\mathsf{nil} : \Pi A{:}\STypeI.\,\mathsf{list}\,A$ and
$\mathsf{cons} : \Pi A{:}\STypeI.\,A \to \mathsf{list}\,A \to
\mathsf{list}\,A$), and a predicate $\mathsf{In} : \Pi
A{:}\STypeI.\,A \to \mathsf{list}\,A \to \SProp$.  The premise
\[
  \varphi \;=\; \Pi A{:}\STypeI.\ \Pi x{:}A.\ \Pi
  l{:}\mathsf{list}\,A.\ \mathsf{In}\,A\,x\,(\mathsf{cons}\,A\,x\,l)
\]
translates to
\[
  \forall A.\ \hastype(A, \lift{\STypeI}) \to
  \forall x.\ \hastype(x, A) \to
  \forall l.\ \hastype(l, \widehat{\mathsf{list}} \app A) \to
  \PR\bigl(\widehat{\mathsf{In}} \app A \app x \app
     (\widehat{\mathsf{cons}} \app A \app x \app l)\bigr).
\]
Every literal carries the type variable $A$; the inversion axiom for
$\mathsf{list}$ is likewise generic in $A$.  \Cref{sec:mono} is
about replacing such formulas by instances at the ground types at
which they are used, e.g.\ $A \coloneqq \mathsf{nat}$;
\cref{sec:specialization} then turns $\hastype(l,
\widehat{\mathsf{list}} \app \widehat{\mathsf{nat}})$ into
$\guard{\mathsf{list\,nat}}(l)$ and flattens the applications; and
\cref{sec:erasure} deletes the guards that monotonicity renders
redundant.
\end{example}
